% ============================================================================
% nb11_beklenti.tex -- 11_encoder_rgb_mixed.ipynb'den ne bekliyoruz?
% Bağımsız derlenir:  latexmk -xelatex nb11_beklenti.tex
% ============================================================================
\documentclass[11pt, paper=a4, DIV=14, parskip=half, numbers=noenddot]{scrartcl}

\usepackage{fontspec}
\usepackage{polyglossia}
\setmainlanguage{turkish}
% Dosya adıyla yükleme: fontlar TinyTeX'in texmf ağacında, macOS CoreText
% onları aile adıyla göremez -- kpathsea dosya adıyla bulur.
\setmainfont{texgyrepagella}[Extension=.otf, UprightFont=*-regular,
  BoldFont=*-bold, ItalicFont=*-italic, BoldItalicFont=*-bolditalic]
\setsansfont{texgyreheros}[Extension=.otf, UprightFont=*-regular,
  BoldFont=*-bold, ItalicFont=*-italic, BoldItalicFont=*-bolditalic]
\setmonofont{texgyrecursor}[Extension=.otf, UprightFont=*-regular,
  BoldFont=*-bold, ItalicFont=*-italic, BoldItalicFont=*-bolditalic,
  Scale=0.92]
\setlength{\emergencystretch}{3em}

\usepackage{xcolor}
\definecolor{anarenk}{HTML}{1B3A6B}
\definecolor{vurgu}{HTML}{C8551A}
\definecolor{yesilvurgu}{HTML}{1E7A52}
\definecolor{grimetin}{HTML}{4A4A46}
\definecolor{acikgri}{HTML}{EDEBE4}
\definecolor{kutucerceve}{HTML}{B9B6A9}
% Excalidraw havası: pastel dolgular, yumuşak köşeler
\definecolor{pastelmavi}{HTML}{DCE9F8}
\definecolor{pastelturuncu}{HTML}{FBE4CE}
\definecolor{pastelyesil}{HTML}{DDF0E2}
\definecolor{pastelgri}{HTML}{EFEDE6}

\addtokomafont{section}{\color{anarenk}\sffamily}
\addtokomafont{subsection}{\color{anarenk}\sffamily}
\addtokomafont{disposition}{\rmfamily}

\usepackage{booktabs}
\usepackage{array}
\usepackage{caption}
\captionsetup{font=small, labelfont={sf,bf,color=anarenk}, margin=1em}
\usepackage{float}
\usepackage{tikz}
\usetikzlibrary{arrows.meta, positioning, calc}

\tikzset{
  kutu/.style={draw=#1!60!grimetin, fill=#1, rounded corners=7pt,
               line width=1.1pt, align=center, font=\sffamily\small,
               inner sep=6pt, minimum height=2.5em},
  kutu/.default=pastelgri,
  donuk/.style={kutu=pastelgri, draw=kutucerceve, text=grimetin,
                dash pattern=on 5pt off 3pt},
  ok/.style={-{Stealth[length=3mm]}, line width=1.1pt, color=grimetin},
  etiket/.style={font=\sffamily\footnotesize, color=grimetin, align=center},
}

\newcommand{\code}[1]{\texttt{\small #1}}

\begin{document}

\begin{center}
  {\LARGE\sffamily\bfseries\color{anarenk}
   Notebook 11 --- RGB Windows in the Batch}\\[3pt]
  {\sffamily Beklenti ve ölçüm planı: \code{11\_encoder\_rgb\_mixed.ipynb}
   \quad$\cdot$\quad Ağustos 2026}
\end{center}
\vspace{2pt}

\section{Amaç ve hipotez}

EdgeTAM'ın \emph{image encoder}'ı SA-1B ile, yani yer seviyesinden çekilmiş
RGB fotoğraflarla ön-eğitildi. Bu projenin şimdiye kadarki eğitimi
(\code{07}) encoder'ı termale uyarlıyor --- bir \emph{modality gap}
kapanıyor. Ama ikinci bir boşluk daha var ve modaliteden bağımsız:
\textbf{yüksek irtifadan dik bakış ve 20--30 piksellik hedefler}, SA-1B'de
her iki modalitede de yok. Bu bir \emph{domain gap} ve RGB'de de termalde de
aynı şekilde açık.

Notebook 11, \code{07}'nin birebir kopyası: aynı schedule, aynı seed, aynı
held-out uçuşlar. \textbf{Tek değişken veri} --- batch'lere iki RGB kaynağı
eklenir. Böylece 07 ile 11'in final sayıları arasındaki her fark yalnızca RGB
pencerelere atfedilebilir.

\begin{itemize}\itemsep2pt
  \item \textbf{H1 (hedef):} Havadan RGB ile eğitim, küçük RGB hedeflerde
        hazır (stock) EdgeTAM'ı geçer --- encoder RGB'yi zaten biliyor,
        öğrenmesi gereken \emph{domain}.
  \item \textbf{H2 (sigorta):} Termal pencereler batch'in çoğunluğu olmaya
        devam ettiği için termal skor düşmez; RGB, drift'e karşı
        regularization gibi davranır.
  \item \textbf{H3 (bonus, belirsiz):} Paylaşılan feature'lar üzerinden RGB
        çeşitliliği termal küçük-hedef skorunu da bir miktar yukarı çekebilir.
\end{itemize}

\section{Mimari: ne eğitiliyor, ne donuk}

\begin{figure}[H]
\centering
\begin{tikzpicture}[node distance=5mm and 8mm]
  % --- Lane 1: stage A -----------------------------------------------------
  \node[etiket, anchor=west] (sa) at (0, 0)
        {\textbf{Stage A} --- feature distillation, etiketsiz
         (hizalı RGB--termal çiftler)};
  \node[kutu=pastelmavi, anchor=west] (teacher) at (0, -1.15)
        {SAM 2.1 Hiera-B+\\\footnotesize teacher (frozen), RGB girdi};
  \node[kutu=pastelyesil, right=16mm of teacher] (stuA)
        {Image Encoder\\\footnotesize student, termal girdi};
  \draw[ok] (teacher) -- node[etiket, above=0.5mm] {cosine\\feature loss}
        (stuA);

  % --- Lane 2: stage B -----------------------------------------------------
  \node[etiket, anchor=west] (sb) at (0, -3.4)
        {\textbf{Stage B} --- supervised fine-tune, karışık batch (bu
         notebook)};
  \node[kutu=pastelturuncu, anchor=west] (win) at (0, -4.65)
        {512$\times$512\\window};
  \node[kutu=pastelyesil, right=of win] (enc)
        {Image Encoder\\\footnotesize RepViT + FPN};
  \node[kutu=pastelyesil, right=of enc] (dec)
        {Prompt Enc.\\+ Mask Decoder};
  \node[kutu=pastelturuncu, right=of dec] (out) {mask logits\\+ IoU};
  \draw[ok] (win) -- (enc);
  \draw[ok] (enc) -- (dec);
  \draw[ok] (dec) -- (out);
  \draw[ok, dashed, color=yesilvurgu]
       (stuA.south) to[bend left=20]
       node[etiket, right=2mm, pos=0.45] {başlangıç ağırlıkları} (enc.north);

  % --- Lane 3: frozen memory path + legend ---------------------------------
  \node[donuk, anchor=west, text width=92mm] (mem) at (0, -6.6)
        {Memory Attention $\cdot$ Memory Encoder $\cdot$ Spatial Perceiver\\
         \textbf{her aşamada donuk} (video = stage C)};
  \node[kutu=pastelyesil, right=5mm of mem, minimum height=1.6em,
        inner sep=3pt] (l1) {\footnotesize eğitilir};
  \node[donuk, right=3mm of l1, minimum height=1.6em, inner sep=3pt] (l2)
        {\footnotesize donuk};
\end{tikzpicture}
\caption{İki aşama, tek encoder. Stage A yalnızca hizalı (RGB, termal)
çiftler okur, hiç etiket görmez; stage B'de encoder head ile birlikte, onda
bir öğrenme hızıyla eğitilir. Memory yolu hiçbir aşamada hareket etmez ---
modality shift bir \emph{encoder} problemidir. Teacher bu kolda SAM 2.1'dir;
DINOv3 ayrı bir deney kolu olarak \code{10}'dadır (bkz.~\S\ref{sec:kapsam}).}
\end{figure}

\section{Veri karışımı}

Beş kaynak, tek batch havuzu. Karışım oranını ayrı bir hiperparametre
yapmadık: her kaynak, pencere havuzuna kare sayısı oranında katılır.

\begin{table}[H]
\centering\small
\begin{tabular}{@{}llllll@{}}
\toprule
Kaynak & Modality & Çözünürlük & Etiket & Rol & Katkı \\
\midrule
VTUAV VIS      & termal & 1920$\times$1080 & instance mask & eğitim + \textbf{skor} & güvenilir GT \\
VTUAV VIS      & RGB    & 1920$\times$1080 & instance mask & eğitim + \textbf{skor} & ölçülebilir RGB \\
SegFly         & termal & 640$\times$512   & semantic map  & yalnız eğitim & hacim, 3 irtifa \\
SegFly dilimi  & RGB    & 4000$\times$3000 & semantic map  & yalnız eğitim & \emph{küçük hedef} \\
Kust4K         & termal & 640$\times$512   & semantic map  & yalnız eğitim & en ince etiket \\
\bottomrule
\end{tabular}
\caption{\code{DATASETS} listesi. Semantic map'ten \emph{reconstruct} edilen
instance'lar (connected components) hedef olarak gürültülü olabileceğinden
asla skorlanmaz; test sayısı yalnızca elle çizilmiş VTUAV maskelerinden
gelir.}
\end{table}

Üç not:

\begin{itemize}\itemsep2pt
  \item \textbf{Kust4K'nın \%29'u dışarıda --- otomatik.} Set, 4\,024
        karesinin 1\,160'ında bir modaliteyi kasıtlı olarak bozar (sensor
        failure simülasyonu) ve bozuk kareleri beş manifest dosyasında
        listeler. İndirme hücresi manifestleri her zaman çeker; okuyucu
        (\code{SPECS["kust4k"].exclude}) bu kareleri daha indeksleme
        aşamasında düşürür. Elle yapılacak bir şey yok.
  \item \textbf{SegFly RGB tam değil, dilim.} RGB yarısı 20\,606 kare /
        156\,GiB'dir; tamamı PNG'ye açılsa $\sim$570\,GB eder. Bunun yerine
        eşit aralıklı shard'lardan 3\,000 kare alınır (her sahne ve her
        irtifa temsil edilir), yayıncının JPEG baytları aynen yazılır:
        $\sim$23\,GB. 4000$\times$3000'den kesilen 512'lik pencerede bir
        araç gerçekten $\sim$20 piksel --- deployment'ın yaşadığı rejim.
  \item \textbf{VTUAV RGB bedava.} Kareler ve maskeleri (\code{rgb/},
        \code{mask/rgb/}) termal koşunun zaten indirdiği üç arşivin içinde.
\end{itemize}

\section{Sızıntısız split ve iki referans}

Split kareler arasında değil \textbf{uçuşlar arasında} kesilir ve bir uçuşun
termal ve RGB kareleri \emph{aynı tarafa} düşer (unit test ile sabitlendi).
Sonuç: held-out bir uçuş hiçbir modalitede görülmemiştir ve 11'in termal test
pencereleri 07'ninkilerle \textbf{birebir aynıdır}.

\begin{figure}[H]
\centering
\begin{tikzpicture}[node distance=4mm and 6mm]
  \node[kutu=pastelmavi]  (f1) {uçuş \code{car\_08}\\\footnotesize termal + RGB};
  \node[kutu=pastelmavi, right=of f1]  (f2) {uçuş \code{bike\_09}\\\footnotesize termal + RGB};
  \node[kutu=pastelmavi, right=of f2]  (f3) {uçuş \code{ped\_15}\\\footnotesize termal + RGB};
  \node[kutu=pastelyesil, below=9mm of f1, minimum width=27mm] (tr) {train};
  \node[kutu=pastelturuncu, below=9mm of f2, minimum width=27mm] (va) {val};
  \node[kutu=pastelgri, below=9mm of f3, minimum width=27mm] (te) {test};
  \draw[ok] (f1) -- (tr);
  \draw[ok] (f2) -- (va);
  \draw[ok] (f3) -- (te);
  \node[etiket, below=2mm of va]
        {bir uçuşun \emph{iki modalitesi birlikte} taşınır --- çapraz modalite sızıntısı yok};
\end{tikzpicture}
\caption{Sequence-level split. Aynı uçuşun RGB'si train'e, termali teste
düşemez; skor ezber değil genelleme ölçer.}
\end{figure}

Eval artık her sayıyı modalite başına ayrı raporlar; böylece iki soru iki
ayrı satırdan okunur:

\begin{table}[H]
\centering\small
\begin{tabular}{@{}llll@{}}
\toprule
Metrik & Referans & Beklenti & Karar kuralı \\
\midrule
\code{rgb / small IoU} ($<$32\,px) & stock EdgeTAM & \textcolor{yesilvurgu}{belirgin artış} &
  \textbf{H1 testi} --- ana hedef \\
\code{rgb / mean IoU} & stock EdgeTAM & artış & ikincil \\
\code{thermal / small IoU} & notebook 07 & $\Delta \geq -0.01$ &
  düşerse: 07 kalır, RGB $\rightarrow$ LoRA \\
\code{thermal / mean IoU} & notebook 07 & değişmez & kontrol \\
\bottomrule
\end{tabular}
\caption{Karar tablosu. ``Small'' kovası kasıtlı olarak ana metriktir: havadan
bir setin ortalama IoU'sunu büyük araçlar domine eder, deployment'ın derdi
ise 20 pikselden küçük hedeflerdir. Termal guard aşılırsa sonuç ``RGB termali
sulandırdı'' demektir; o durumda çözüm iki ayrı model değil, ortak taban +
modalite başına $\sim$0.3M parametrelik LoRA adapter'dır.}
\end{table}

\section{Maliyet ve çalıştırma}

Notebook uçtan uca tek \emph{Run all}'dur: veri indirme, doğrulama
(md5), indeksleme, stage A, üç eğitim koşusu (fine-tune, LoRA,
distilled+fine-tune) ve stock dahil dört checkpoint'in test skoru. Gated
model yoktur; SAM 2.1 teacher anonim iner. Gereken tek onay Drive mount'tur.

\begin{table}[H]
\centering\small
\begin{tabular}{@{}lll@{}}
\toprule
Kalem & 07 & 11 \\
\midrule
İndirme & $\sim$43\,GB (VTUAV) + SegFly + Kust4K & + $\sim$23\,GB (SegFly RGB dilimi) \\
Disk & $\sim$60\,GB & $\sim$85\,GB \\
GPU (A100) & $\sim$20\,dk stage A + 3$\times\sim$40\,dk eğitim & aynı schedule, aynı adım sayısı \\
Ek hesap & --- & yalnızca 12\,MP JPEG decode (loader'da) \\
\bottomrule
\end{tabular}
\caption{11'in tek ek maliyeti veridir; adım sayısı ve schedule 07 ile
aynıdır, GPU süresi kayda değer değişmez. Çıktılar kendi Drive klasörüne
(\code{edgetam-encoder/rgb}) yazılır; 07 ile aynı anda ayrı runtime'da
koşturulabilir.}
\end{table}

\section{Kapsam dışı}
\label{sec:kapsam}

\begin{itemize}\itemsep2pt
  \item \textbf{DINOv3 teacher bu kolda yok --- kasıtlı.} Her notebook tek
        değişken oynatır: \code{10} teacher'ı değiştirir (SAM 2.1
        $\rightarrow$ DINOv3), \code{11} veriyi değiştirir (+RGB). İkisi
        birden değişseydi fark hangisine atfedilecekti? Önce iki kol ayrı
        ölçülür; DINOv3 \code{10}'da kazanırsa kombinasyon (DINOv3 + RGB)
        bir sonraki koşudur.
  \item \textbf{SegFly RGB'nin tamamı} (20\,606 kare) --- dilim yeterli
        sinyal veriyorsa büyütmek bir satırlık değişikliktir
        (\code{rows=3000}).
  \item \textbf{Video ve memory yolu} --- stage C; encoder sayıları
        ölçülmeden açılmaz.
\end{itemize}

\vspace{0.5em}
\noindent\textcolor{grimetin}{\small Tek cümlelik özet: \emph{11, ``termali
bozmadan RGB'yi bedavaya alabilir miyiz?'' sorusunu, 07 ile birebir
karşılaştırılabilir iki sayıya --- \code{rgb/small IoU} vs stock ve
\code{thermal/small IoU} vs 07 --- indirger.}}

\end{document}
